\documentclass[11pt]{article}
\usepackage[margin=0.8in]{geometry}
\usepackage{amsmath,amssymb,mathtools}
\usepackage{booktabs,array,enumitem,microtype}
\setlength{\parindent}{0pt}
\setlength{\parskip}{0.62em}

\newcommand{\ket}[1]{\lvert #1\rangle}
\newcommand{\bra}[1]{\langle #1\rvert}
\newcommand{\I}{\mathbb{I}}
\newcommand{\CNOT}[2]{\operatorname{CNOT}_{#1\to #2}}
\newcommand{\CZ}{\operatorname{CZ}}

\begin{document}

\begin{center}
  {\Large\bfseries Problem 1 -- Question 5}\\[0.3em]
  {\large Gate construction}
\end{center}

\section*{5(a): SWAP from CNOT gates}

Use the chronological sequence
\[
 \boxed{\CNOT{1}{2}\;,\quad\CNOT{2}{1}\;,\quad\CNOT{1}{2}.}
\]
On a computational basis state \(\ket{x,y}\), the bit values evolve as
\[
 (x,y)\mapsto(x,x\oplus y)
 \mapsto(y,x\oplus y)\mapsto(y,x).
\]
Linearity gives the same result for arbitrary states
\(\ket\psi\ket\phi\). Three CNOTs are necessary and sufficient if CNOT
is the only available two-qubit gate.

\section*{5(b): Two distance-two CNOTs on a four-qubit line}

Assume the desired directions are \(q_0\to q_2\) and \(q_1\to q_3\).
The following nearest-neighbor schedule uses six CNOTs:
\[
\begin{array}{c|l}
\text{layer}&\text{gates}\\ \hline
1&\CNOT{0}{1}\\
2&\CNOT{1}{2}\\
3&\CNOT{0}{1}\quad\text{and}\quad\CNOT{2}{3}\\
4&\CNOT{1}{2}\\
5&\CNOT{2}{3}.
\end{array}
\]
The two gates in layer \(3\) act on disjoint pairs and run in parallel.
Tracking the four classical basis bits gives
\[
 (x_0,x_1,x_2,x_3)
 \longmapsto
 (x_0,x_1,x_2\oplus x_0,x_3\oplus x_1),
\]
so the induced unitary is exactly the product of the two requested
CNOTs. The count is \(6\) and the depth is \(5\). A SWAP-routing
construction would use \(8\) CNOTs, so preserving the temporary parity
information saves two gates.

\section*{5(c): \(R_{ZZ}\) from CNOT and \(R_Z\)}

Because
\[
 \CNOT{1}{2}\,(\I\otimes Z)\,\CNOT{1}{2}=Z\otimes Z,
\]
conjugating an \(R_Z\) rotation gives
\[
\boxed{
 R_{ZZ}(\theta)
 =\CNOT{1}{2}
   \bigl(\I\otimes R_Z(\theta)\bigr)
   \CNOT{1}{2}.
}
\]
Chronologically: CNOT \(1\to2\), \(R_Z(\theta)\) on qubit \(2\), then
the same CNOT again. The circuit uses two CNOTs and one \(R_Z\).

\section*{5(d): Controlled phase from \(R_{ZZ}\)}

The projector onto \(\ket{11}\) is
\[
 \ket{11}\!\bra{11}
 =\frac14(\I-Z_1)(\I-Z_2)
 =\frac14(\I-Z_1-Z_2+Z_1Z_2).
\]
Since all four terms commute,
\begin{align*}
\operatorname{CP}(\theta)
 &=\exp\bigl(i\theta\ket{11}\!\bra{11}\bigr)\\
 &=e^{i\theta/4}
   R_{Z,1}\left(\frac{\theta}{2}\right)
   R_{Z,2}\left(\frac{\theta}{2}\right)
   R_{ZZ}\left(-\frac{\theta}{2}\right).
\end{align*}
Hence, up to the physically irrelevant global phase \(e^{i\theta/4}\),
\[
\boxed{
\operatorname{CP}(\theta)
\equiv
R_{Z,1}\left(\frac{\theta}{2}\right)
R_{Z,2}\left(\frac{\theta}{2}\right)
R_{ZZ}\left(-\frac{\theta}{2}\right).
}
\]
The resulting diagonal entries are \(1,1,1,e^{i\theta}\), after the
global phase is restored.

\section*{5(e): Exponential of an arbitrary Pauli string}

Let
\[
 P=P_1\otimes\cdots\otimes P_n,\qquad
 P_j\in\{\I,X,Y,Z\},
\]
and let \(S=\{j:P_j\neq\I\}\) have size \(w\). For each \(j\in S\),
apply a basis change \(B_j\) satisfying \(B_jP_jB_j^\dagger=Z\):
\[
 B_j=
 \begin{cases}
 H,&P_j=X,\\
 HS^\dagger,&P_j=Y,\\
 \I,&P_j=Z.
 \end{cases}
\]
For \(Y\), the chronological order is \(S^\dagger\) and then \(H\).
After these changes, \(P\) has become a product of \(Z\)'s.

Choose one qubit \(t\in S\) as a parity target. Apply CNOTs from every
other \(j\in S\setminus\{t\}\) to \(t\), apply \(R_Z(2\theta)\) to
\(t\), undo the CNOTs in reverse order, and finally undo all \(B_j\).
The middle rotation is
\[
 R_Z(2\theta)=e^{-i\theta Z_t},
\]
and parity conjugation replaces \(Z_t\) by \(\prod_{j\in S}Z_j\).
Thus the complete circuit implements
\[
 \boxed{e^{-i\theta P}.}
\]
It uses \(2(w-1)\) CNOTs, one \(R_Z\), and at most \(2w\) one-qubit
basis-change gates. A balanced parity tree reduces the CNOT depth from
\(2(w-1)\) to \(2\lceil\log_2w\rceil\) when connectivity permits.

\section*{5(f): The specified two-qubit matrix}

Use computational-basis order
\(\ket{00},\ket{01},\ket{10},\ket{11}\), and define
\[
R_Y(\vartheta)=
\begin{pmatrix}
\cos(\vartheta/2)&-\sin(\vartheta/2)\\
\sin(\vartheta/2)&\cos(\vartheta/2)
\end{pmatrix}.
\]
The requested real matrix \(U\) has the exact factorization
\[
\boxed{
\begin{aligned}
U={}&
\CZ\,
\left[\I\otimes R_Y\left(-\frac{13\pi}{12}\right)\right]\,
\CZ\\
&\times
\left[R_Y\left(\frac{3\pi}{4}\right)
\otimes R_Y\left(-\frac{5\pi}{12}\right)\right]\,
\CZ\,
\left[R_Y\left(\frac{3\pi}{4}\right)\otimes\I\right].
\end{aligned}}
\]
Thus the chronological pulse schedule is
\[
\begin{array}{c|l}
1&R_Y(3\pi/4)\text{ on }q_0\\
2&\CZ\\
3&R_Y(3\pi/4)\text{ on }q_0
   \quad\text{and}\quad R_Y(-5\pi/12)\text{ on }q_1\\
4&\CZ\\
5&R_Y(-13\pi/12)\text{ on }q_1\\
6&\CZ .
\end{array}
\]
Direct multiplication yields
\[
\begin{pmatrix}
\frac12&\frac1{\sqrt2}&\frac12&0\\
\frac12&-\frac1{\sqrt2}&\frac12&0\\
\frac1{2\sqrt2}&0&-\frac1{2\sqrt2}&\frac{\sqrt3}{2}\\
\frac{\sqrt3}{2\sqrt2}&0&-\frac{\sqrt3}{2\sqrt2}&-\frac12
\end{pmatrix}.
\]
The two simultaneous rotations in layer \(3\) give a total circuit
depth of \(6\), with three CZ gates and four \(R_Y\) pulses.

\paragraph{Determinant note.}
As printed, the matrix is real orthogonal with determinant \(-1\), not
\(+1\). Therefore it is not literally in \(SU(4)\). Multiplying it by
the global phase \(e^{i\pi/4}\) gives determinant \(1\), because
\(\det(e^{i\pi/4}U)=e^{i\pi}\det U=1\). This global phase has no
observable effect, so the displayed circuit implements the intended
two-qubit operation.

\end{document}
